% Puzzle difficulty: the paper. The prose follows docs/writing.md. Every
% number, table, figure and board is written by `lab paper puzzle-difficulty`
% into generated/, from the model's committed results, and read back with
% \V{key} and \block{name}.
\documentclass[11pt]{article}
% The preamble every paper of the lab shares: % The preamble every paper of the lab shares: % The preamble every paper of the lab shares: \input{../../../papers/preamble}
% from models/<model>/paper/, \input{../preamble} from papers/<name>/.
% Boards come from chessboard and xskak, plots from pgfplots. A paper prints no
% number by hand: `lab paper` writes every number, table and figure under the
% paper's generated/ folder, and paper.tex says what they mean.
\usepackage[LSBC4,T1]{fontenc}
\usepackage[utf8]{inputenc}
\usepackage{lmodern}
\usepackage[margin=1in]{geometry}
\usepackage{microtype}
\usepackage{amsmath}
\usepackage{booktabs}
\usepackage{tabularx}
\usepackage{array}
\usepackage{xcolor}
\usepackage[obeyspaces,spaces]{url}
\usepackage{xskak}
\usepackage{pgfplots}
\usepgfplotslibrary{groupplots}
\usetikzlibrary{positioning,arrows.meta,calc,fit}
\pgfplotsset{compat=1.18}
\usepackage{caption}
\usepackage{subcaption}
\usepackage{enumitem}
\usepackage[section]{placeins}
\usepackage[hidelinks]{hyperref}

% A generated value by its key: generated/values.tex sets them, \V{key} prints
% one, and a key nobody set stops the build.
\makeatletter
\newcommand{\setvalue}[2]{\expandafter\gdef\csname labvalue@#1\endcsname{#2}}
\newcommand{\V}[1]{\@ifundefined{labvalue@#1}%
  {\PackageError{lab}{No generated value named '#1'}{Run `lab paper` for this model again.}}%
  {\csname labvalue@#1\endcsname}}
\makeatother

% Chart colours. A run keeps its colour in every figure of every paper.
\definecolor{labBlue}{HTML}{2A78D6}
\definecolor{labOrange}{HTML}{EB6834}
\definecolor{labAqua}{HTML}{1BAF7A}
\definecolor{labYellow}{HTML}{EDA100}
\definecolor{labMagenta}{HTML}{E87BA4}
\definecolor{labGreen}{HTML}{008300}
\definecolor{labViolet}{HTML}{6250D6}
\definecolor{labGray}{HTML}{898781}
\definecolor{labInk}{HTML}{0B0B0B}
\definecolor{labSecondary}{HTML}{52514E}
\definecolor{labGrid}{HTML}{E1E0D9}
\definecolor{labAxis}{HTML}{C3C2B7}
\definecolor{labWash}{HTML}{CDE2FB}

\pgfplotsset{
  lab/.style={
    font=\sffamily\footnotesize,
    axis line style={labAxis},
    tick style={labAxis},
    tick label style={labSecondary},
    label style={labSecondary},
    title style={labInk, font=\sffamily\small\bfseries, align=left},
    grid style={labGrid, line width=0.3pt},
    major grid style={labGrid, line width=0.3pt},
    axis lines*=left,
    tick align=outside,
    clip=false,
    legend style={draw=none, fill=none, font=\sffamily\footnotesize, text=labSecondary,
      cells={anchor=west}, legend columns=-1, column sep=6pt, at={(0,1.03)}, anchor=south west},
    every axis plot/.append style={line width=1.1pt, line join=round, line cap=round},
    scaled ticks=false,
    /pgf/number format/1000 sep={,},
  },
  lab marks/.style={mark=*, mark size=1.9pt, mark options={draw=white, line width=0.6pt}},
}

\setchessboard{boardfontencoding=LSBC4, boardfontfamily=skaknew, blackfieldcolor=black!18,
  whitefieldcolor=white, setfontcolors, showmover=false, label=true, labelfontsize=6pt,
  boardfontsize=12pt}

% A move or line in UCI, a position in FEN (breakable at a slash), and a run's name.
\newcommand{\uci}[1]{\texttt{#1}}
\newcommand{\fen}[1]{{\footnotesize\path{#1}}}
\newcommand{\run}[1]{\texttt{#1}}
\newcolumntype{L}{>{\raggedright\arraybackslash}X}
\newcolumntype{R}{>{\raggedleft\arraybackslash}X}

\setlength{\parskip}{0.4em}
\setlength{\parindent}{0pt}
\captionsetup{font=small, labelfont=bf}

% from models/<model>/paper/, \input{../preamble} from papers/<name>/.
% Boards come from chessboard and xskak, plots from pgfplots. A paper prints no
% number by hand: `lab paper` writes every number, table and figure under the
% paper's generated/ folder, and paper.tex says what they mean.
\usepackage[LSBC4,T1]{fontenc}
\usepackage[utf8]{inputenc}
\usepackage{lmodern}
\usepackage[margin=1in]{geometry}
\usepackage{microtype}
\usepackage{amsmath}
\usepackage{booktabs}
\usepackage{tabularx}
\usepackage{array}
\usepackage{xcolor}
\usepackage[obeyspaces,spaces]{url}
\usepackage{xskak}
\usepackage{pgfplots}
\usepgfplotslibrary{groupplots}
\usetikzlibrary{positioning,arrows.meta,calc,fit}
\pgfplotsset{compat=1.18}
\usepackage{caption}
\usepackage{subcaption}
\usepackage{enumitem}
\usepackage[section]{placeins}
\usepackage[hidelinks]{hyperref}

% A generated value by its key: generated/values.tex sets them, \V{key} prints
% one, and a key nobody set stops the build.
\makeatletter
\newcommand{\setvalue}[2]{\expandafter\gdef\csname labvalue@#1\endcsname{#2}}
\newcommand{\V}[1]{\@ifundefined{labvalue@#1}%
  {\PackageError{lab}{No generated value named '#1'}{Run `lab paper` for this model again.}}%
  {\csname labvalue@#1\endcsname}}
\makeatother

% Chart colours. A run keeps its colour in every figure of every paper.
\definecolor{labBlue}{HTML}{2A78D6}
\definecolor{labOrange}{HTML}{EB6834}
\definecolor{labAqua}{HTML}{1BAF7A}
\definecolor{labYellow}{HTML}{EDA100}
\definecolor{labMagenta}{HTML}{E87BA4}
\definecolor{labGreen}{HTML}{008300}
\definecolor{labViolet}{HTML}{6250D6}
\definecolor{labGray}{HTML}{898781}
\definecolor{labInk}{HTML}{0B0B0B}
\definecolor{labSecondary}{HTML}{52514E}
\definecolor{labGrid}{HTML}{E1E0D9}
\definecolor{labAxis}{HTML}{C3C2B7}
\definecolor{labWash}{HTML}{CDE2FB}

\pgfplotsset{
  lab/.style={
    font=\sffamily\footnotesize,
    axis line style={labAxis},
    tick style={labAxis},
    tick label style={labSecondary},
    label style={labSecondary},
    title style={labInk, font=\sffamily\small\bfseries, align=left},
    grid style={labGrid, line width=0.3pt},
    major grid style={labGrid, line width=0.3pt},
    axis lines*=left,
    tick align=outside,
    clip=false,
    legend style={draw=none, fill=none, font=\sffamily\footnotesize, text=labSecondary,
      cells={anchor=west}, legend columns=-1, column sep=6pt, at={(0,1.03)}, anchor=south west},
    every axis plot/.append style={line width=1.1pt, line join=round, line cap=round},
    scaled ticks=false,
    /pgf/number format/1000 sep={,},
  },
  lab marks/.style={mark=*, mark size=1.9pt, mark options={draw=white, line width=0.6pt}},
}

\setchessboard{boardfontencoding=LSBC4, boardfontfamily=skaknew, blackfieldcolor=black!18,
  whitefieldcolor=white, setfontcolors, showmover=false, label=true, labelfontsize=6pt,
  boardfontsize=12pt}

% A move or line in UCI, a position in FEN (breakable at a slash), and a run's name.
\newcommand{\uci}[1]{\texttt{#1}}
\newcommand{\fen}[1]{{\footnotesize\path{#1}}}
\newcommand{\run}[1]{\texttt{#1}}
\newcolumntype{L}{>{\raggedright\arraybackslash}X}
\newcolumntype{R}{>{\raggedleft\arraybackslash}X}

\setlength{\parskip}{0.4em}
\setlength{\parindent}{0pt}
\captionsetup{font=small, labelfont=bf}

% from models/<model>/paper/, \input{../preamble} from papers/<name>/.
% Boards come from chessboard and xskak, plots from pgfplots. A paper prints no
% number by hand: `lab paper` writes every number, table and figure under the
% paper's generated/ folder, and paper.tex says what they mean.
\usepackage[LSBC4,T1]{fontenc}
\usepackage[utf8]{inputenc}
\usepackage{lmodern}
\usepackage[margin=1in]{geometry}
\usepackage{microtype}
\usepackage{amsmath}
\usepackage{booktabs}
\usepackage{tabularx}
\usepackage{array}
\usepackage{xcolor}
\usepackage[obeyspaces,spaces]{url}
\usepackage{xskak}
\usepackage{pgfplots}
\usepgfplotslibrary{groupplots}
\usetikzlibrary{positioning,arrows.meta,calc,fit}
\pgfplotsset{compat=1.18}
\usepackage{caption}
\usepackage{subcaption}
\usepackage{enumitem}
\usepackage[section]{placeins}
\usepackage[hidelinks]{hyperref}

% A generated value by its key: generated/values.tex sets them, \V{key} prints
% one, and a key nobody set stops the build.
\makeatletter
\newcommand{\setvalue}[2]{\expandafter\gdef\csname labvalue@#1\endcsname{#2}}
\newcommand{\V}[1]{\@ifundefined{labvalue@#1}%
  {\PackageError{lab}{No generated value named '#1'}{Run `lab paper` for this model again.}}%
  {\csname labvalue@#1\endcsname}}
\makeatother

% Chart colours. A run keeps its colour in every figure of every paper.
\definecolor{labBlue}{HTML}{2A78D6}
\definecolor{labOrange}{HTML}{EB6834}
\definecolor{labAqua}{HTML}{1BAF7A}
\definecolor{labYellow}{HTML}{EDA100}
\definecolor{labMagenta}{HTML}{E87BA4}
\definecolor{labGreen}{HTML}{008300}
\definecolor{labViolet}{HTML}{6250D6}
\definecolor{labGray}{HTML}{898781}
\definecolor{labInk}{HTML}{0B0B0B}
\definecolor{labSecondary}{HTML}{52514E}
\definecolor{labGrid}{HTML}{E1E0D9}
\definecolor{labAxis}{HTML}{C3C2B7}
\definecolor{labWash}{HTML}{CDE2FB}

\pgfplotsset{
  lab/.style={
    font=\sffamily\footnotesize,
    axis line style={labAxis},
    tick style={labAxis},
    tick label style={labSecondary},
    label style={labSecondary},
    title style={labInk, font=\sffamily\small\bfseries, align=left},
    grid style={labGrid, line width=0.3pt},
    major grid style={labGrid, line width=0.3pt},
    axis lines*=left,
    tick align=outside,
    clip=false,
    legend style={draw=none, fill=none, font=\sffamily\footnotesize, text=labSecondary,
      cells={anchor=west}, legend columns=-1, column sep=6pt, at={(0,1.03)}, anchor=south west},
    every axis plot/.append style={line width=1.1pt, line join=round, line cap=round},
    scaled ticks=false,
    /pgf/number format/1000 sep={,},
  },
  lab marks/.style={mark=*, mark size=1.9pt, mark options={draw=white, line width=0.6pt}},
}

\setchessboard{boardfontencoding=LSBC4, boardfontfamily=skaknew, blackfieldcolor=black!18,
  whitefieldcolor=white, setfontcolors, showmover=false, label=true, labelfontsize=6pt,
  boardfontsize=12pt}

% A move or line in UCI, a position in FEN (breakable at a slash), and a run's name.
\newcommand{\uci}[1]{\texttt{#1}}
\newcommand{\fen}[1]{{\footnotesize\path{#1}}}
\newcommand{\run}[1]{\texttt{#1}}
\newcolumntype{L}{>{\raggedright\arraybackslash}X}
\newcolumntype{R}{>{\raggedleft\arraybackslash}X}

\setlength{\parskip}{0.4em}
\setlength{\parindent}{0pt}
\captionsetup{font=small, labelfont=bf}

% Written by `lab paper puzzle-difficulty`; do not edit by hand.
\looktrue
\chancefalse
\pairsfalse
\variantsfalse
\gamesfalse
\featuresfalse
\changesfalse
\setvalue{bands.best-ahead.of}{four of the five bands}
\setvalue{bands.count}{five}
\setvalue{bands.lead.narrowest}{$-$4.15}
\setvalue{bands.lead.narrowest.name}{1400-1799}
\setvalue{bands.lead.widest}{+65.6}
\setvalue{bands.lead.widest.name}{2200+}
\setvalue{best.beside.closer}{57}
\setvalue{best.beside.correlation}{0.94}
\setvalue{best.close.median}{157.3}
\setvalue{best.close.p90}{449.1}
\setvalue{best.close.within.100}{34}
\setvalue{best.close.within.200}{60}
\setvalue{best.close.within.300}{76}
\setvalue{best.device}{NVIDIA A100-SXM4-80GB}
\setvalue{best.epochs.run}{12}
\setvalue{best.gain.of-lift}{12}
\setvalue{best.gain.percent}{8.4}
\setvalue{best.mae}{204.8}
\setvalue{best.parameters}{6{,}563{,}649}
\setvalue{best.r2}{0.7112}
\setvalue{best.reliability.groups}{20}
\setvalue{best.reliability.highest.happened}{2276}
\setvalue{best.reliability.highest.said}{2316}
\setvalue{best.reliability.lowest.happened}{650}
\setvalue{best.reliability.lowest.said}{618}
\setvalue{best.reliability.spread}{738}
\setvalue{best.reliability.widest-gap}{40.5}
\setvalue{best.rows-a-second}{15{,}100}
\setvalue{best.score}{270.2}
\setvalue{best.score.high}{271.6}
\setvalue{best.score.low}{268.8}
\setvalue{best.seconds}{13 min}
\setvalue{best.spearman}{0.8484}
\setvalue{challenger.name}{\run{transformer-medium-hand}}
\setvalue{code.commit}{\texttt{85d0c35}}
\setvalue{context.lift.share}{36}
\setvalue{cost.seconds}{35 min}
\setvalue{cost.seconds.networks}{33 min}
\setvalue{cost.seconds.others}{2 min}
\setvalue{data.fingerprint}{\texttt{d14731fa4425}}
\setvalue{data.licence}{CC0 1.0}
\setvalue{data.name}{\texttt{puzzles-v1}}
\setvalue{data.rows}{1{,}127{,}915}
\setvalue{data.rows.test}{112{,}450}
\setvalue{data.share.test}{10}
\setvalue{data.share.train}{85}
\setvalue{data.share.valid}{5}
\setvalue{examples.count}{three}
\setvalue{floor.constant}{1462.4}
\setvalue{floor.score}{502.9}
\setvalue{grid.smallest-cell}{30}
\setvalue{label.mean}{1462.4}
\setvalue{label.q5}{693.0}
\setvalue{label.q95}{2338.0}
\setvalue{label.sd}{503.3}
\setvalue{ladder.at-last-epoch.of}{three of the four networks}
\setvalue{look.date}{22 November 2026}
\setvalue{look.resamples}{200}
\setvalue{look.smallest-compared}{200}
\setvalue{look.test-rows}{112{,}450}
\setvalue{manifest.kept.percent}{18}
\setvalue{manifest.source-rows-read}{6{,}193{,}530}
\setvalue{method.gbm.features.percent}{80}
\setvalue{method.gbm.leaves}{127}
\setvalue{method.gbm.patience}{100}
\setvalue{method.gbm.rate}{0.05}
\setvalue{method.gbm.rounds}{4{,}000}
\setvalue{method.gbm.rows.percent}{80}
\setvalue{method.resamples}{1{,}000}
\setvalue{networks.ahead.of}{all four networks}
\setvalue{networks.patience}{three}
\setvalue{networks.warmup.percent}{5}
\setvalue{recipe.filters.max-rating-deviation}{100}
\setvalue{recipe.filters.min-plays}{100}
\setvalue{recipe.sample.percent}{25}
\setvalue{reference.features}{38}
\setvalue{reference.leans.1.feature}{quiet moves in the line}
\setvalue{reference.leans.1.share}{22}
\setvalue{reference.leans.2.feature}{line ends in mate}
\setvalue{reference.leans.2.share}{11}
\setvalue{reference.leans.3.feature}{pieces the line moves}
\setvalue{reference.leans.3.share}{10}
\setvalue{reference.lift.percent}{41}
\setvalue{reference.score}{294.9}
\setvalue{reference.trees}{2{,}970}
\setvalue{results.date}{22 November 2026}
\setvalue{run.conv-small.gain}{+4.105}
\setvalue{run.gbm-context.name}{\run{gbm-context}}
\setvalue{run.gbm-context.score}{428.8}
\setvalue{run.transformer-medium-hand.gain}{+24.68}
\setvalue{run.transformer-medium-hand.over-twin}{+13.8}
\setvalue{run.transformer-medium.gain}{+10.91}
\setvalue{run.transformer-medium.parameters.millions}{6.6}
\setvalue{run.transformer-small.gain}{+1.847}
\setvalue{runs.count}{seven}
\setvalue{slice.backward.no.happened}{1452}
\setvalue{slice.backward.yes.happened}{1557}
\setvalue{slice.band.-1000.challenger.off}{+101}
\setvalue{slice.band.-1000.reference.off}{+165}
\setvalue{slice.band.2200+.challenger.off}{$-$341}
\setvalue{slice.band.2200+.reference.off}{$-$424}
\setvalue{slice.capture.no.happened}{1459}
\setvalue{slice.capture.yes.happened}{1466}
\setvalue{slice.check.no.happened}{1724}
\setvalue{slice.check.yes.happened}{1300}
\setvalue{slice.deviation.-75.gain}{+1.170}
\setvalue{slice.deviation.-75.happened}{1669}
\setvalue{slice.deviation.90+.gain}{+54.34}
\setvalue{slice.deviation.90+.happened}{1226}
\setvalue{slice.first-piece.king.happened}{1750}
\setvalue{slice.first-piece.pawn.happened}{1674}
\setvalue{slice.first-piece.queen.happened}{1357}
\setvalue{slice.mate.no.gain}{+26.49}
\setvalue{slice.mate.no.happened}{1613}
\setvalue{slice.mate.yes.gain}{+20.03}
\setvalue{slice.mate.yes.happened}{1119}
\setvalue{slice.moves.1.challenger.off}{$-$26}
\setvalue{slice.moves.1.gain}{+12.52}
\setvalue{slice.moves.2.gain}{+23.39}
\setvalue{slice.moves.3.gain}{+30.18}
\setvalue{slice.moves.ahead.of}{all five groups}
\setvalue{slice.moves.behind}{zero}
\setvalue{slice.plays.-200.gain}{+29.29}
\setvalue{slice.plays.5000+.gain}{+12.40}
\setvalue{slice.quiet.0.gain}{+22.17}
\setvalue{slice.quiet.0.happened}{1301}
\setvalue{slice.quiet.1.happened}{1759}
\setvalue{slice.quiet.2.gain}{+35.43}
\setvalue{slice.quiet.2.happened}{1872}
\setvalue{slice.sacrifice.no.happened}{1443}
\setvalue{slice.sacrifice.yes.happened}{1580}
\setvalue{slice.themes.ahead}{30}
\setvalue{slice.themes.behind}{zero}
\setvalue{slice.themes.compared}{30}
\setvalue{slice.themes.gain.least}{+12.49}
\setvalue{slice.themes.gain.least.group}{mateIn1}
\setvalue{slice.themes.gain.most}{+39.47}
\setvalue{slice.themes.gain.most.group}{quietMove}
\setvalue{task.context.count}{five}
\setvalue{task.hand.count}{34}
\setvalue{task.marks.count}{14}
\setvalue{task.reference}{\run{gbm-hand}}


\begin{document}

\papertitle{Predicting the Difficulty of Chess Puzzles from the Position and
  Its Solution}
  {Results as of \V{results.date}\quad\textperiodcentered\quad dataset \V{data.name},
   build \V{data.fingerprint}\quad\textperiodcentered\quad code \V{code.commit}}

\begin{abstract}
The rating of a Lichess puzzle measures how difficult the puzzle is for human
solvers: a solver with the same rating is expected to solve it about half of
the time. We train models that predict this rating from a puzzle's position
and its solution, so that a tactic without a rating, such as one missed in a
player's own game, can be placed on the same scale. The dataset contains
\V{data.rows} puzzles with stable ratings, split by the game each puzzle was
extracted from. We compare a constant baseline, gradient-boosted decision
trees (GBDT) on hand-crafted features, and convolutional and transformer
networks that read the board as 64 square tokens. The best of
\V{runs.count} models is a transformer that also receives the hand-crafted
features and whose encoder is pretrained on engine evaluations. On
\V{data.rows.test} test puzzles its root-mean-square error (RMSE) is
\V{best.score} rating points (95\% CI \V{best.score.low} to
\V{best.score.high}), compared with \V{reference.score} for the GBDT baseline
and \V{floor.score} for the constant baseline. Compared with the GBDT
baseline, \V{networks.ahead.of} are significantly better. The best model
ranks puzzles well, with a Spearman correlation of \V{best.spearman}, but its
mean absolute error is \V{best.mae} rating points. Its predictions are
therefore suitable for ordering tactics by difficulty and not yet for
reporting the rating of an individual tactic.
\end{abstract}

\section{Introduction}
\label{sec:intro}

A chess engine identifies the best move in a position. It does not indicate
how difficult that move is for a person to find. A one-move fork that wins a
queen and a quiet rook manoeuvre that wins the same queen five moves later
receive similar evaluations, yet most players find only the first.

Lichess publishes several million puzzles, each extracted from a real game at
a moment when one side made a mistake~\cite{lichessdb}. Every puzzle has a
Glicko-2 rating~\cite{glicko2}: each attempt is scored as a game between the
solver and the puzzle, which the solver wins by playing the complete
solution. Once a puzzle has been attempted often enough its rating is stable,
and it is the best available measure of the puzzle's difficulty for human
solvers.

A model that predicts this rating from a position and its solution can rate
any tactic, including the tactics in a player's own games, which no other
player has attempted. This work supports a chess coaching application that
would use such a model in two ways: in a game review, to distinguish a missed
tactic that was within the player's reach from one well above the player's
level, and to select training positions of suitable difficulty.

We evaluate a fixed sequence of models on \V{data.rows} puzzles: a constant
baseline, GBDT on scalar context features, GBDT with hand-crafted features of
the position and the solution, and neural networks that read the board
directly, with and without the hand-crafted features. Because the sequence is
fixed, each step measures the contribution of one component. Our main
findings are the following.
\begin{itemize}
\item \textbf{Hand-crafted features of the solution carry most of the
signal.} The GBDT baseline has an RMSE of \V{reference.score} rating points,
\V{reference.lift.percent}\% below the constant baseline (\V{floor.score}).
The context features alone, of which the solution length is the informative
one, account for \V{context.lift.share}\% of this
improvement\iffeatures, and removing the features that describe the solution
line costs the baseline more than removing any other group
(Section~\ref{sec:feature-groups})\fi.
\item \textbf{Neural networks improve on GBDT.} Compared with the GBDT
baseline, \V{networks.ahead.of} are significantly better. The best model has
a gain of \V{best.gain} RMSE points (95\% CI \V{best.gain.low} to
\V{best.gain.high}), which is \V{best.gain.percent}\% of the baseline's error
(Section~\ref{sec:results}).
\ifchangestart\ifchangeattacks
\item \textbf{Initialisation and attention bias matter more than model
size.} Pretraining the encoder on engine evaluations and adding an
attack-aware attention bias each improve the medium transformer by more than
doubling its training epochs does. A larger transformer and a
convolution-transformer hybrid are significantly worse than the medium
transformer (Section~\ref{sec:ablations}).
\fi\fi
\ifmore
\item \textbf{More training data helps GBDT more than the transformer.} With
four times the training puzzles and a fixed number of optimisation steps, the
gain is \V{pair.gbm-hand-quadruple.gbm-hand.gain} RMSE points for GBDT and
\V{pair.transformer-medium-hand-quadruple.transformer-medium-hand.gain} for
the transformer (Section~\ref{sec:training-size}).
\fi
\item \textbf{The best model is accurate on average and imprecise for a
single puzzle.} Its mean absolute error is \V{best.mae} rating
points\iflook, and \V{best.close.within.200}\% of its predictions fall within
200 points of the true rating. Its mean prediction is within
\V{best.reliability.widest-gap} points of the mean true rating at every
predicted level (Section~\ref{sec:error})\fi.
\end{itemize}

\section{Related Work}
\label{sec:related}

\paragraph{Puzzle difficulty.}
Predicting the Lichess rating of a puzzle from its position and solution was
the task of data-mining competitions at IEEE BigData 2024 and FedCSIS
2025~\cite{puzzlecup}. The winning entry of the 2025 competition combines a
fine-tuned Maia-2~\cite{maia2}, a model of human moves, with hand-crafted and
engine-derived features~\cite{puzzlecupwinner}. Mi{\l}osz and Kapusta~\cite{glickformer} introduce
GlickFormer, a transformer that encodes the sequence of positions along the
solution and predicts the puzzle's Glicko-2 rating. Our study differs in its
aim. We do not combine all available sources of information to maximise
accuracy. We start from the board alone and add hand-crafted features,
pretraining and architectural changes one at a time, to measure the
contribution of each. We also use our own sample of the puzzle database and
our own test set, so our scores are not comparable with the competition
results.

\paragraph{Models of human play.}
Maia~\cite{maia} and Maia-2~\cite{maia2} predict the moves that players of a
given rating make. Anderson et al.~\cite{humanerror} find that the difficulty
of a position predicts human error better than the player's skill or the time
available. These studies model behaviour in games. Puzzle ratings offer a
complementary signal: a direct measurement, from many attempts, of how hard
one specific tactic is.

\section{Task and Data}
\label{sec:data}

\subsection{Task}

The input is the position that the solver sees, the opponent's preceding
move, and the solution line. The target is the puzzle's Lichess rating.
Table~\ref{tab:brief} summarises the task.

\begin{table}[htbp]\centering\small
\begin{tabularx}{\textwidth}{lL}
\toprule
Input & The position presented to the solver, the opponent's preceding move
and the solution line. The origin and destination squares of the moves of the
line are marked in auxiliary input planes. \\
Target & The puzzle's Lichess rating, in rating points. \\
Data & The Lichess puzzle database (\V{data.licence}): a fixed
\V{recipe.sample.percent}\% sample of the puzzles with a stable rating. \\
Split & By the game a puzzle was extracted from: \V{data.share.train}\%
training, \V{data.share.valid}\% validation, \V{data.share.test}\% test. \\
Primary metric & RMSE on the test set, in rating points (lower is better). \\
\bottomrule
\end{tabularx}
\caption{Summary of the task.}
\label{tab:brief}
\end{table}

\subsection{Dataset}

The source is the Lichess puzzle database~\cite{lichessdb}. Each record gives
the puzzle's identifier, the position before the opponent's move, the moves
of the solution with the opponent's move first, the puzzle's rating and
rating deviation, the number of attempts, the theme tags, and the game the
puzzle was extracted from.

We keep the puzzles whose rating is stable: a Glicko-2 rating deviation of at
most \V{recipe.filters.max-rating-deviation} after at least
\V{recipe.filters.min-plays} attempts. Of these we keep a fixed
\V{recipe.sample.percent}\% sample, selected by a hash of the puzzle
identifier, so that the dataset can be rebuilt exactly and a larger sample
contains this one. The resulting dataset has \V{data.rows} puzzles
(Table~\ref{tab:data}).

The split is determined by a hash of the source game, so no game contributes
puzzles to more than one of the training, validation and test sets. The same
grouping is used for the bootstrap (Section~\ref{sec:protocol}). The dataset
is not stored in the repository. Its recipe is, and a build of the dataset is
identified by a fingerprint, a hash of its contents. All models in this paper
were trained and evaluated on build \V{data.fingerprint}\ifmore; the models
of Section~\ref{sec:training-size} that use more training puzzles were
trained on a larger build and evaluated on the validation and test sets of
this one\fi. The evaluation code refuses to rank models evaluated on
different builds.

\begin{table}[htbp]\centering\small
\block{data}
\caption{Dataset statistics by split.}
\label{tab:data}
\end{table}

\iflook
The source contains \V{manifest.source-rows-read} puzzles, of which
\V{manifest.kept.percent}\% remain after filtering and sampling. The ratings
have a mean of \V{label.mean} and a standard deviation of \V{label.sd}, and
90\% of them lie between \V{label.q5} and \V{label.q95}
(Figure~\ref{fig:label}, Table~\ref{tab:describe}). The standard deviation is
also the RMSE of a model that predicts the mean rating for every puzzle.

\begin{figure}[htbp]\centering
\block{label}
\caption{Distribution of the puzzle ratings in the dataset.}
\label{fig:label}
\end{figure}

\begin{table}[htbp]\centering\small
\block{describe}
\caption{Summary statistics of selected variables over the whole dataset:
the mean, five quantiles, the minimum and the maximum.}
\label{tab:describe}
\end{table}
\fi

\section{Methods}
\label{sec:methods}

\subsection{Input representation}
\label{sec:input}

% Input representation: identical for every task, so each model paper inputs
% this file under its own heading and then adds what is specific to its task.
% The prose follows docs/writing.md.
All tasks in this project share one input representation, so two tasks differ
in their target and in the information a model is given, and not in how a
position is encoded.

\paragraph{Square tokens.}
A position is encoded as 64 tokens, one per square. Token 0 denotes an empty
square, tokens 1 to 6 the pawn, knight, bishop, rook, queen and king of the
side to move, which we call the \emph{mover}, and tokens 7 to 12 the
opponent's pieces. When Black is to move, the board is flipped vertically and
the colours are swapped, so that the mover always plays up the board. A
position with White to move and the corresponding colour-reversed position
with Black to move therefore have the same encoding, and a model does not
have to learn the two cases separately.

\paragraph{Auxiliary input planes.}
Information about individual squares is supplied as binary planes over the 64
squares, which a network receives together with the tokens. This task uses
\V{task.marks.count} planes (Table~\ref{tab:marks}). The first marks the en
passant target square, and the others mark the origin and destination squares
of moves. A plane identifies the squares of a move and carries no other
information about it, such as whether the move is a capture or a check.

\begin{table}[htbp]\centering\small
\block{marks}
\caption{Auxiliary input planes. Squares are given from the mover's point of
view, after the board has been flipped.}
\label{tab:marks}
\end{table}

\ifsequence
\paragraph{Sequences of positions.}
An example of this task is not a single position. It is the sequence of one
player's moves in a game, and each element of the sequence is a position
encoded as above: the board before the move, with the opponent's preceding
move and the move played marked in the auxiliary planes. Scalar inputs
therefore come at two levels, one for each move and one for the game.

\paragraph{Per-move features.}
Each position of the sequence has \V{task.steps.count} scalar features of its
own (Table~\ref{tab:steps}): the four castling rights and the task's
information about that move. A network receives them together with the
position they belong to.

\begin{table}[htbp]\centering\small
\block{steps}
\caption{Per-move features: the scalar inputs attached to each position of a
sequence.}
\label{tab:steps}
\end{table}

\paragraph{Scalar context features.}
Information about the game as a whole is given as \V{task.context.count}
scalar context features (Table~\ref{tab:reads}). Every model receives them,
including the gradient-boosted trees, which never see a board. For a network,
both kinds of scalar feature are standardised with the mean and standard
deviation of the training set.

\begin{table}[htbp]\centering\small
\block{reads}
\caption{Scalar context features: the game-level inputs given to every
model.}
\label{tab:reads}
\end{table}

\paragraph{Hand-crafted features.}
Gradient-boosted trees cannot take a sequence of boards as input, so we also
compute \V{task.hand.count} aggregate features over the player's moves
(Table~\ref{tab:hand} in the appendix): how much win probability the moves
lost on average, at worst and how often beyond fixed thresholds, how the clock
was used, and how many moves were available. These features have two roles.
They are the inputs of the GBDT baseline, and so define the level that a
network reading every position has to exceed. They can also be given to a
network as additional context features, which tests whether the network
already extracts the same information from the sequence.

\else
\paragraph{Scalar context features.}
Information that does not belong to a square is given as scalar context
features: the four castling rights and the scalar inputs specific to the
task. This task uses \V{task.context.count} (Table~\ref{tab:reads}). Every
model receives them, including the gradient-boosted trees, which never see
the board. For a network they are standardised with the mean and standard
deviation of the training set.

\begin{table}[htbp]\centering\small
\block{reads}
\caption{Scalar context features: the inputs given to every model in addition
to the board.}
\label{tab:reads}
\end{table}

\paragraph{Hand-crafted features.}
Gradient-boosted trees cannot take 64 tokens as input, so we also compute
\V{task.hand.count} hand-crafted features of the position
(Table~\ref{tab:hand} in the appendix): material, mobility, available checks
and captures, hanging pieces, passed pawns and attacks near each king, and,
where the task gives a line of moves, features of that line. These features
have two roles. They are the inputs of the GBDT baseline, and so define the
level that a network reading the board has to exceed. They can also be given
to a network as additional context features, which tests whether the network
already extracts the same information from the board.

\fi

\iflook
\begin{figure}[htbp]\centering
\block{encoding}
\ifsequence
\caption{One position of one example, shown twice: the position before the
player's most costly move of the game. Left: as played, with arrows for the
opponent's preceding move (grey) and the player's move (blue). Right: as
encoded, oriented so that the mover plays up the board, with the squares set
in an auxiliary plane shaded grey for the opponent's move and blue for the
player's move. Every position of the sequence is encoded in this way.}
\else
\caption{One example from the dataset, shown twice. Left: as played, with
arrows for the opponent's preceding move (grey) and the move the example
concerns (blue). Right: as encoded, oriented so that the mover plays up the
board, with the squares set in an auxiliary plane shaded grey for the
opponent's move and, where the task marks a line of moves, blue for the
mover's moves and orange for the replies.}
\fi
\label{fig:encoding}
\end{figure}
\fi


\paragraph{Encoding of the solution.}
The task is to rate a known tactic, so every model is given the solution. The
auxiliary planes mark the opponent's preceding move, the first three moves of
the solver and the two replies between them, each by its origin and
destination square, and one further plane marks every square that a later
move of the line leaves or reaches. The only scalar context feature besides
the castling rights is the number of moves the solver has to find. A network
therefore sees where the pieces move. It is not told whether a move is a
capture, a check or a sacrifice, and has to infer this from the board. The
hand-crafted features state exactly these properties.

\subsection{Models and training}
\label{sec:models}

% Models and training: the same sequence of models for every task, so each
% model paper inputs this file under its own heading.
% The prose follows docs/writing.md.
Every task is given the same sequence of models, so that each step of the
sequence measures the contribution of one component. A model is fitted on the
training set, selected on the validation set and scored once on the test set.
Table~\ref{tab:runs} lists the models of this task.

\paragraph{Constant baseline.}
\run{constant} predicts the same value for every example: the mean of the
training targets (for a binary target, the base rate), or their median when
the primary metric is an absolute error. A useful model has to improve on it.

\paragraph{Gradient-boosted decision trees.}
\ifcontextgbdt
\run{gbm-context} is a gradient-boosted decision tree model (GBDT), trained
with LightGBM~\cite{lightgbm} on the scalar context features alone. It never
sees the board, and so measures how much of the target the context already
determines. \run{gbm-hand} adds the hand-crafted features.
\else
\run{gbm-hand} is a gradient-boosted decision tree model (GBDT), trained with
LightGBM~\cite{lightgbm} on the scalar context features and the hand-crafted
features.
\fi
We refer to \run{gbm-hand} as \emph{the GBDT baseline}: every other model is
compared with it on the same test examples. GBDT is trained for at most
\V{method.gbm.rounds} boosting rounds with up to \V{method.gbm.leaves} leaves
per tree and a learning rate of \V{method.gbm.rate}. Each tree is fitted on a
random \V{method.gbm.rows.percent}\% of the examples and
\V{method.gbm.features.percent}\% of the features, and boosting stops when the
validation score has not improved for \V{method.gbm.patience} rounds.

\paragraph{Neural networks.}
A network consists of an encoder, which maps the square tokens, the auxiliary
planes and the context features to a vector, and a prediction head, a
two-layer perceptron that maps this vector to the output. We use two encoder
architectures, so that the results show how much the architecture matters for
the task.

The \emph{transformer} encoder~\cite{attention} treats each square as a token.
The embedding of a token is the sum of a learned embedding of its piece, a
learned embedding of its square and a linear projection of the square's
auxiliary planes. The context features are projected by a small perceptron
and enter as one additional token. Every token attends to every other token
in every layer, so distant squares interact from the first layer on. To give
attention the geometry of the board, each attention head adds a learned bias
that depends on the file and rank offset between two
squares~\cite{relative}: files, ranks, diagonals and knight moves all
correspond to fixed offsets. The layers use pre-layer normalisation, and the
encoder outputs the context token concatenated with the mean of the 64 square
tokens.

The \emph{residual convolutional network} (ResNet) is a stack of residual
blocks of $3\times3$ convolutions over the $8\times8$ board, the architecture
that AlphaZero introduced for chess~\cite{alphazero}. Its input is one binary
plane per token value together with the auxiliary planes. A linear projection
of the context features is added at every square after the first convolution,
and the encoder outputs the mean and the maximum over the squares of the
final feature map. Unlike attention, a convolution relates only neighbouring
squares, so the receptive field grows with depth.

\ifsequence
\paragraph{Sequence model.}
An example of this task is a sequence of positions, so a network has two
stages below its prediction head. One of the two encoders above is applied to
every position of the sequence, with the position's per-move features in
place of the context features, and produces one vector per position. A
two-layer transformer then operates on these vectors in the order in which
the moves were played: each vector is a token, a learned positional embedding
is added to it, and the game-level context features enter as one additional
token. The prediction head receives that token concatenated with the mean of
the position tokens. Sequences in a batch are padded to the length of the
longest, and padded elements are masked: they are not encoded and no token
attends to them, so a prediction does not depend on the other examples in
the batch.

Model sizes are hyperparameters. \run{conv-small} is small enough to train
on a CPU. An example consists of tens of positions, so
\run{transformer-small} and \run{transformer-medium} are trained on a GPU.
\run{transformer-medium-hand} is \run{transformer-medium} with the
hand-crafted features added to its context features.
Table~\ref{tab:settings} in the appendix lists the hyperparameters.
\else
Model sizes are hyperparameters. \run{conv-small} and \run{transformer-small}
are small enough to train on a CPU. \run{transformer-medium} is trained on a
GPU, and \run{transformer-medium-hand} is the same network with the
hand-crafted features added to its context features.
Table~\ref{tab:settings} in the appendix lists the hyperparameters.
\fi

\paragraph{Training.}
All networks are trained in the same way: with AdamW~\cite{adamw}, a learning
rate that increases linearly over the first \V{networks.warmup.percent}\% of
the optimisation steps and then follows a cosine schedule to zero, gradient
clipping, and mixed precision on a GPU. The loss matches the primary metric:
mean squared error for a target scored by RMSE, a Huber loss for a target
scored by MAE, and binary cross-entropy for a probability. A regression
target is standardised with the mean and standard deviation of the training
set. After every epoch the network is evaluated on the validation set. We
keep the weights of the best epoch, and stop training when
\V{networks.patience} consecutive epochs have not improved on it.
\ifmore
A model trained on a larger training set is given proportionally fewer
epochs, and a correspondingly shorter patience
(Section~\ref{sec:ablations}).
\fi

\begin{table}[htbp]\centering\footnotesize
\block{runs}
\caption{The models, best first. \emph{Selected} is the number of trees or
the training epoch chosen on the validation set. \emph{Time} is the total
wall-clock time for training and evaluation.}
\label{tab:runs}
\end{table}


\section{Evaluation Protocol}
\label{sec:protocol}

\paragraph{Metrics.}
The primary metric, fixed before the first full training run, is the RMSE on
the test set in rating points; lower is better. We report three secondary
metrics: the mean absolute error (MAE), the coefficient of determination
(R\textsuperscript{2}), and Spearman's rank correlation, which measures only whether the
puzzles are ordered correctly. We also report results separately for
\V{bands.count} bands of the true rating, because an error of a given size
matters most for puzzles near a player's own level.

% Evaluation protocol: the same for every task. Each model paper names its
% own metrics under its section heading and inputs the rest from here.
% The prose follows docs/writing.md.
\paragraph{Confidence intervals.}
Examples from the same game are not independent, so we resample whole games.
The test games are drawn with replacement \V{method.resamples} times, the
metric is computed each time on the examples of the games drawn, and the
2.5th and 97.5th percentiles of these values form the 95\% confidence
interval (CI) that we report with every score.

\paragraph{Paired comparisons.}
To compare two models we compute both scores on the same resampled games and
take their difference. We call this difference the \emph{gain} of the first
model over the second, and sign it so that a positive gain means that the
first model is better. A gain is \emph{statistically significant} when its
95\% CI excludes zero, and is then set in bold in the tables. Unless stated
otherwise, the second model is the GBDT baseline. Two models tend to do well
and badly on the same games, so a paired gain can be significant even when
the two models' separate intervals overlap. These intervals reflect the
sampling of test games only. Each network was trained with a single random
seed, and the variation between training runs is not included.

\paragraph{Model selection.}
\ifchance
The number of trees, the training epoch whose weights are kept and the
calibration parameters are chosen on the validation set.
\else
The number of trees and the training epoch whose weights are kept are chosen
on the validation set.
\fi
The test set is used once per model, for the scores reported here.

\ifchance
\paragraph{Calibration.}
Before a model is scored, its predicted probabilities are rescaled by Platt
scaling~\cite{platt}: a logistic function of the log-odds whose slope and
intercept are fitted on the validation set. The transformation is monotone,
so it leaves a model's ranking of the examples unchanged and affects only its
log loss and its calibration. What remains of a model's miscalibration after
this step cannot be removed by a linear correction of its log-odds.
\fi

\ifpairs
\paragraph{Acceptance criterion.}
Section~\ref{sec:ablations} evaluates modifications of the models. A
modification is \emph{accepted} only if the modified model is better than the
unmodified one on the validation set and its paired gain on the test set is
statistically significant. A modification that is not accepted is still
reported, as a measurement of its effect.
\fi

\paragraph{Tables and figures.}
A model has the same colour and marker in every figure, and the GBDT baseline
is always drawn in orange. The values plotted in a figure are also given in a
table, in the text or in the appendix. All numbers, tables and figures are
generated from the committed result files by the command that builds the
paper.


\section{Results}
\label{sec:results}

\subsection{Main results}

Table~\ref{tab:leaderboard} and Figure~\ref{fig:scores} give the main
results. The constant baseline predicts \V{floor.constant} for every puzzle
and has an RMSE of \V{floor.score}, the standard deviation of the ratings.
GBDT on the context features alone, that is on the solution length and the
castling rights, reaches \V{run.gbm-context.score} and thereby achieves
\V{context.lift.share}\% of the improvement that the GBDT baseline achieves
over the constant baseline. The GBDT baseline, with the hand-crafted features
of the position and the solution, reaches \V{reference.score}, which is
\V{reference.lift.percent}\% below the constant baseline.

\begin{table}[htbp]\centering\small
\block{leaderboard}
\caption{Main results on the \V{data.rows.test} test puzzles, best first.
The gain is the paired difference from \V{task.reference} on the same
resampled games. It is positive when the model is better, and bold when
statistically significant.}
\label{tab:leaderboard}
\end{table}

\begin{figure}[htbp]\centering
\block{scores}
\caption{Test RMSE of every model with its 95\% CI. The vertical line marks
the GBDT baseline.}
\label{fig:scores}
\end{figure}

Compared with the GBDT baseline, \V{networks.ahead.of} are significantly
better (Figure~\ref{fig:gains}). Within the fixed sequence of models, accuracy
increases with model size and with the hand-crafted features. The small
transformer and the small ResNet, with
\V{run.transformer-small.parameters.millions} and
\V{run.conv-small.parameters.millions} million parameters, improve on the
GBDT baseline only slightly (gains of \V{run.transformer-small.gain} and
\V{run.conv-small.gain} RMSE points). The medium transformer, with
\V{run.transformer-medium.parameters.millions} million parameters, gains
\V{run.transformer-medium.gain}, and the same transformer with hand-crafted
features gains \V{run.transformer-medium-hand.gain}. At these sizes the
number of parameters matters more than the choice between a convolutional
and a transformer encoder.

\begin{figure}[htbp]\centering
\block{gains}
\caption{Paired gain in RMSE of each network over the GBDT baseline, with
its 95\% CI. A model whose interval lies entirely to the right of the line is
significantly better than the baseline.}
\label{fig:gains}
\end{figure}

The hand-crafted features remain useful to a network.
\ifpairs
Compared with \run{transformer-medium}, \run{transformer-medium-hand} is
\V{pair.transformer-medium-hand.transformer-medium.verdict} on the test set
(gain \V{pair.transformer-medium-hand.transformer-medium.gain}; 95\% CI
\V{pair.transformer-medium-hand.transformer-medium.gain.low} to
\V{pair.transformer-medium-hand.transformer-medium.gain.high}).
\else
The test scores of the medium transformer with and without them differ by
\V{run.transformer-medium-hand.over-twin} RMSE points, positive when the
features help.
\fi
With the standard schedule, \V{run.transformer-medium.depth} transformer
layers do not extract from the board all the information that the
hand-crafted features provide.
\ifvariants
Section~\ref{sec:ablations} shows that most of this difference disappears
when both models are trained for twice as long.
\fi

The best model's gain is \V{best.gain.percent}\% of the GBDT baseline's
error, and \V{best.gain.of-lift}\% of the improvement of the GBDT baseline
over the constant baseline. GBDT on hand-crafted features therefore accounts
for most of what is predicted, and the network adds a smaller improvement.

\subsection{Ranking accuracy and absolute error}

The best model orders puzzles well and is less accurate in absolute terms.
Its Spearman correlation with the true rating is \V{best.spearman} and its
R\textsuperscript{2} is \V{best.r2} (Table~\ref{tab:metrics}), while its MAE is
\V{best.mae} rating points. Chess players are conventionally grouped into
rating classes 200 points wide, so the typical error is close to one class.
The model can therefore order the tactics of a game by difficulty, which is
what a game review needs in order to say that one missed tactic was routine
and another was hard. It cannot yet state the rating of a single tactic to
within one class.

\begin{table}[htbp]\centering\small
\block{metrics}
\caption{All metrics for every model on the test set. RMSE and MAE are in
rating points (lower is better); for R\textsuperscript{2} and Spearman's rank correlation
higher is better. The constant baseline has no rank correlation.}
\label{tab:metrics}
\end{table}

\subsection{Results by rating band}

Table~\ref{tab:bands} gives the RMSE by band of the true rating. Because the
bands are defined by the target, they favour predictions near the mean: the
constant baseline has its smallest error in the middle band and its largest
errors in the extreme bands. Errors should therefore be compared between
models within a band, and not between bands.

\begin{table}[htbp]\centering\small
\block{bands-table}
\caption{Test RMSE by band of the true puzzle rating. The bands are defined
by the target, so values are comparable within a column and not between
columns.}
\label{tab:bands}
\end{table}

The best model has a lower RMSE than the GBDT baseline in
\V{bands.best-ahead.of}. The difference is largest for puzzles rated
\V{bands.lead.widest.name} (\V{bands.lead.widest} points) and smallest for
puzzles rated \V{bands.lead.narrowest.name} (\V{bands.lead.narrowest}
points). The network's advantage is thus concentrated at the two extremes, on
the easiest and the hardest puzzles (Figure~\ref{fig:bands}).

\begin{figure}[htbp]\centering
\block{bands}
\caption{Difference in RMSE between the GBDT baseline and each network, by
rating band (the differences of the values in Table~\ref{tab:bands}, without
intervals). Above the dashed line the network is better.}
\label{fig:bands}
\end{figure}

\subsection{Training curves}

Figure~\ref{fig:curves} shows the validation RMSE of each network of the
fixed sequence after every epoch. Of these, \V{ladder.at-last-epoch.of}
reached their best validation score in the final epoch, so their training
ended while they were still improving.
\ifvariants
A longer schedule may therefore improve them, and
Section~\ref{sec:ablations} reports the same models with a longer schedule.
\else
A longer schedule may therefore improve them, which was not tested.
\fi

\begin{figure}[htbp]\centering
\block{curves}
\caption{Validation RMSE of each network of the fixed sequence after every
epoch. The dashed line is the GBDT baseline's validation RMSE.}
\label{fig:curves}
\end{figure}

\subsection{Feature importance of the GBDT baseline}

The GBDT baseline consists of \V{reference.trees} trees over
\V{reference.features} features. Figure~\ref{fig:leans} shows each feature's
share of the total split gain. The three most important features are
\V{reference.leans.1.feature} (\V{reference.leans.1.share}\% of the gain),
\V{reference.leans.2.feature} (\V{reference.leans.2.share}\%) and
\V{reference.leans.3.feature} (\V{reference.leans.3.share}\%). All three
describe the solution and not the position. Split gain shows which features
the model used. It does not identify what makes a puzzle difficult, because
correlated features divide the credit between them arbitrarily.

\begin{figure}[htbp]\centering
\block{leans}
\caption{The ten most important features of the GBDT baseline, by share of
the total split gain.}
\label{fig:leans}
\end{figure}

\ifchanges
\section{Ablations}
\label{sec:ablations}

% Ablations: what a modification is, the table of paired comparisons and the
% two figures of validation curves. The same for every task, so a model paper
% inputs this file at the start of its ablation section and then reports its
% own findings. Each part appears once the task has the runs it needs, and the
% list below needs at least one of its items.
% The prose follows docs/writing.md.
\ifpairs
The main results compare complete models. An ablation isolates one factor: it
compares a model with its \emph{base}, a model that differs from it in a
single respect. Table~\ref{tab:pairs} lists every such comparison for this
task, with the gain of the modified model over its base on the validation set
and its paired gain on the test set, computed on the same resampled test
games as in the main results. The acceptance criterion of
Section~\ref{sec:protocol} is applied to these pairs. The modifications are
the following.
\begin{itemize}
\ifchangehand
\item \textbf{Hand-crafted features} (suffix \code{-hand}): the network
receives the hand-crafted features as additional context features.
\fi
\ifchangeepochs
\item \textbf{Longer schedule} (suffix \code{-long}): twice the training
epochs, with the same warm-up fraction and cosine schedule.
\fi
\ifchangestart
\item \textbf{Pretraining} (suffix \code{-pretrained}): the transformer
encoder is initialised from a transformer of the same size that was trained
to predict the engine's win probability for positions from the Lichess
evaluation database~\cite{lichessdb}. The piece and square embeddings and all
transformer layers are copied. The prediction head is initialised randomly,
as are the projections of the auxiliary planes that the pretraining task does
not have and, unless the two tasks have the same context features, the
projection of the context features. Training then follows the standard
schedule.
\fi
\ifchangeattacks
\item \textbf{Attack-aware attention bias} (suffix \code{-attacks}): the
relative position bias encodes where two squares lie and not what stands
between them. With this modification each attention head learns four
additional biases for a pair of squares, applied when the piece on one square
attacks the other, when it is attacked from the other, and for the same two
relations through exactly one intervening piece, as in a pin or a skewer. The
relations are computed from the tokens without search, and the biases are
initialised to zero.
\fi
\ifchangegeometry
\item \textbf{No relative position bias} (suffix \code{-no-geometry}): the
relative position bias is removed. This configuration measures the
contribution of the bias and is not a candidate model.
\fi
\ifchangeconv
\item \textbf{ResNet encoder} (prefix \code{conv-medium}): a ResNet with
squeeze-and-excitation in place of the transformer, trained at the higher
peak learning rate used for the convolutional models.
\fi
\ifchangehybrid
\item \textbf{Hybrid encoder} (prefix \code{hybrid-medium}): a few residual
convolutional blocks followed by transformer layers over the resulting
square representations.
\fi
\ifchangelarger
\item \textbf{Larger transformer} (prefix \code{transformer-large}): more layers and
a larger width, with the training schedule unchanged.
\fi
\ifgames
\item \textbf{Training-set size}: the same model trained on a fraction or a
multiple of the training games (Section~\ref{sec:training-size}).
\fi
\iffeatures
\item \textbf{Feature groups}: a model trained without one group of
hand-crafted features at a time (Section~\ref{sec:feature-groups}).
\fi
\end{itemize}

\longblock{pairs}{Paired comparisons. Each row compares a model with a
baseline, in most rows the model it modifies. \emph{Change} states how the
two differ. The gain is that of the model over the baseline and is positive
when the model is better. On
the validation set it is the difference between the two scores. On the test
set it is the paired gain, with its 95\% CI, and is set in bold when
statistically significant.}{tab:pairs}
\fi

\ifvariants
Figure~\ref{fig:variants} shows the validation curves of each model of the
fixed sequence whose training schedule or initialisation was modified,
together with the modified models. A modified model has the colour of its
base and a broken line: long dashes for a longer schedule and dots for a
pretrained encoder.

\begin{figure}[htbp]\centering
\block{variants}
\caption{Validation curves of a model of the fixed sequence (solid line) and
of the models that modify its training schedule or initialisation (broken
lines), after every epoch. One panel per base model, on a common scale. The
dashed horizontal line is the GBDT baseline's validation score.}
\label{fig:variants}
\end{figure}
\fi

\ifshapes
Figure~\ref{fig:shapes} shows the validation curves of the models that
change the architecture and nothing about the training data, each in its own
panel beside its base.

\begin{figure}[htbp]\centering
\block{shapes}
\caption{Validation curves of a model with a modified architecture (broken
line) and of its base (solid line), after every epoch. One panel per modified
model, on a common scale. The dashed horizontal line is the GBDT baseline's
validation score.}
\label{fig:shapes}
\end{figure}
\fi


\ifchangeepochs
\subsection{Training schedule}

Doubling the number of epochs, from \V{run.transformer-medium.epochs} to
\V{run.transformer-medium-long.epochs}, improves both medium transformers.
Without hand-crafted features, \run{transformer-medium-long} is
\V{pair.transformer-medium-long.transformer-medium.valid.verdict} than its
base on the validation set and
\V{pair.transformer-medium-long.transformer-medium.verdict} on the test set
(gain \V{pair.transformer-medium-long.transformer-medium.gain} RMSE points;
95\% CI \V{pair.transformer-medium-long.transformer-medium.gain.low} to
\V{pair.transformer-medium-long.transformer-medium.gain.high}), so the
modification is \V{pair.transformer-medium-long.transformer-medium.kept}.
With hand-crafted features, \run{transformer-medium-hand-long} is
\V{pair.transformer-medium-hand-long.transformer-medium-hand.valid.verdict}
than its base on the validation set and
\V{pair.transformer-medium-hand-long.transformer-medium-hand.verdict} on the
test set (gain
\V{pair.transformer-medium-hand-long.transformer-medium-hand.gain}; 95\% CI
\V{pair.transformer-medium-hand-long.transformer-medium-hand.gain.low} to
\V{pair.transformer-medium-hand-long.transformer-medium-hand.gain.high}),
and the modification is
\V{pair.transformer-medium-hand-long.transformer-medium-hand.kept}.

The longer schedule is worth more than the hand-crafted features. It improves
the transformer without hand-crafted features by more than adding the
features does at the standard schedule (a gain of
\V{pair.transformer-medium-hand.transformer-medium.gain}; the two intervals
in Table~\ref{tab:pairs} do not overlap). As a result,
\run{transformer-medium-long} is more accurate than
\run{transformer-medium-hand}: their gains over the GBDT baseline are
\V{run.transformer-medium-long.gain} and \V{run.transformer-medium-hand.gain}.

With the longer schedule the hand-crafted features contribute little.
Compared with \run{transformer-medium-long},
\run{transformer-medium-hand-long} is
\V{pair.transformer-medium-hand-long.transformer-medium-long.verdict} on the
test set, with a gain of
\V{pair.transformer-medium-hand-long.transformer-medium-long.gain} (95\% CI
\V{pair.transformer-medium-hand-long.transformer-medium-long.gain.low} to
\V{pair.transformer-medium-hand-long.transformer-medium-long.gain.high}),
where the corresponding gain at the standard schedule is
\V{pair.transformer-medium-hand.transformer-medium.gain}. The two models
rank puzzles equally well, with Spearman correlations of
\V{run.transformer-medium-hand-long.spearman} and
\V{run.transformer-medium-long.spearman}. The model with hand-crafted
features has the lower RMSE in the highest rating band
(\V{run.transformer-medium-hand-long.band.5} against
\V{run.transformer-medium-long.band.5}) and the higher RMSE in the lowest
(\V{run.transformer-medium-hand-long.band.1} against
\V{run.transformer-medium-long.band.1}; Table~\ref{tab:bands}, without
intervals). We hypothesise that, given enough optimisation steps, the
transformer learns from the board most of what the hand-crafted features
encode, and that the features remain useful for the hardest puzzles.

Neither model used its full schedule. \run{transformer-medium-long} had its
best validation score at epoch \V{run.transformer-medium-long.best-epoch} of
\V{run.transformer-medium-long.epochs} and
\run{transformer-medium-hand-long} at epoch
\V{run.transformer-medium-hand-long.best-epoch}, and each was stopped
\V{networks.patience} epochs later by early stopping
(Figure~\ref{fig:variants}). This is consistent with overfitting to the
training puzzles, so we do not expect a still longer schedule from random
initialisation to help.
\fi

\ifchangestart
\subsection{Pretraining}

A pretrained encoder improves both medium transformers at the standard
schedule. Without hand-crafted features,
\run{transformer-medium-pretrained} is
\V{pair.transformer-medium-pretrained.transformer-medium.valid.verdict} than
its base on the validation set and
\V{pair.transformer-medium-pretrained.transformer-medium.verdict} on the
test set (gain \V{pair.transformer-medium-pretrained.transformer-medium.gain}
RMSE points; 95\% CI
\V{pair.transformer-medium-pretrained.transformer-medium.gain.low} to
\V{pair.transformer-medium-pretrained.transformer-medium.gain.high}), and
the modification is
\V{pair.transformer-medium-pretrained.transformer-medium.kept}. With
hand-crafted features, \run{transformer-medium-hand-pretrained} is
\V{pair.transformer-medium-hand-pretrained.transformer-medium-hand.valid.verdict}
than its base on the validation set and
\V{pair.transformer-medium-hand-pretrained.transformer-medium-hand.verdict}
on the test set (gain
\V{pair.transformer-medium-hand-pretrained.transformer-medium-hand.gain};
95\% CI
\V{pair.transformer-medium-hand-pretrained.transformer-medium-hand.gain.low}
to
\V{pair.transformer-medium-hand-pretrained.transformer-medium-hand.gain.high}),
and the modification is
\V{pair.transformer-medium-hand-pretrained.transformer-medium-hand.kept}.
The best model of this task is \V{best.name}.

\ifchangeepochs
Pretraining gains more than the longer schedule, in half the number of
epochs: \V{pair.transformer-medium-pretrained.transformer-medium.gain}
against \V{pair.transformer-medium-long.transformer-medium.gain} without
hand-crafted features, and
\V{pair.transformer-medium-hand-pretrained.transformer-medium-hand.gain}
against \V{pair.transformer-medium-hand-long.transformer-medium-hand.gain}
with them. In neither case do the two intervals in Table~\ref{tab:pairs}
overlap. At the standard schedule, the pretrained transformer without
hand-crafted features (RMSE \V{run.transformer-medium-pretrained.score}) is
as accurate as the transformer with hand-crafted features trained for twice
as long (\V{run.transformer-medium-hand-long.score}).

With a pretrained encoder the hand-crafted features contribute again. The
pretrained model with hand-crafted features has a test RMSE of
\V{run.transformer-medium-hand-pretrained.score} and the pretrained model
without them one of \V{run.transformer-medium-pretrained.score}, a
difference of \V{run.transformer-medium-hand-pretrained.over-twin} points.
The corresponding paired gains are
\V{pair.transformer-medium-hand.transformer-medium.gain} at the standard
schedule and
\V{pair.transformer-medium-hand-long.transformer-medium-long.gain} with the
longer schedule. The two pretrained models were not compared directly on
resampled games, but the intervals of their gains over the GBDT baseline do
not overlap (Table~\ref{tab:leaderboard}). A plausible explanation is that an
encoder pretrained on engine evaluations represents which side stands better,
which is not what makes a tactic hard to find, whereas the hand-crafted
features describe the moves of the solution directly.
\fi

The pretrained model with hand-crafted features had its best validation
score at epoch \V{run.transformer-medium-hand-pretrained.best-epoch} of
\V{run.transformer-medium-hand-pretrained.epochs}, whereas the pretrained
model without them had it at epoch
\V{run.transformer-medium-pretrained.best-epoch}
(Figure~\ref{fig:variants}). The former was therefore stopped by its schedule
and not by early stopping, and may improve with further training.
\fi

\ifshapes
\subsection{Architecture}

The architecture ablations keep the training data, the prediction head and
the number of epochs of a medium transformer and change the encoder. All of
these models are trained from random initialisation, as their bases are.
\begin{itemize}
\ifchangeconv
\item \textbf{ResNet.} \run{conv-medium-hand} has
\V{run.conv-medium-hand.parameters.millions} million parameters, compared
with \V{run.transformer-medium-hand.parameters.millions} million for the
transformer, and is trained with a peak learning rate of
\V{run.conv-medium-hand.lr} instead of \V{run.transformer-medium-hand.lr}.
Compared with \run{transformer-medium-hand} it is
\V{pair.conv-medium-hand.transformer-medium-hand.verdict} on the test set
(gain \V{pair.conv-medium-hand.transformer-medium-hand.gain}; 95\% CI
\V{pair.conv-medium-hand.transformer-medium-hand.gain.low} to
\V{pair.conv-medium-hand.transformer-medium-hand.gain.high}), and the
modification is \V{pair.conv-medium-hand.transformer-medium-hand.kept}.
\fi
\ifchangehybrid
\item \textbf{Hybrid.} \run{hybrid-medium-hand} has
\V{run.hybrid-medium-hand.convs} residual convolutional blocks followed by
\V{run.hybrid-medium-hand.depth} transformer layers, for a total of
\V{run.hybrid-medium-hand.parameters.millions} million parameters. Compared
with \run{transformer-medium-hand} it is
\V{pair.hybrid-medium-hand.transformer-medium-hand.verdict} on the test set
(gain \V{pair.hybrid-medium-hand.transformer-medium-hand.gain}; 95\% CI
\V{pair.hybrid-medium-hand.transformer-medium-hand.gain.low} to
\V{pair.hybrid-medium-hand.transformer-medium-hand.gain.high}), and the
modification is \V{pair.hybrid-medium-hand.transformer-medium-hand.kept}.
\fi
\ifchangelarger
\item \textbf{Larger transformer.} \run{transformer-large-hand} has
\V{run.transformer-large-hand.depth} layers of width
\V{run.transformer-large-hand.dim} instead of
\V{run.transformer-medium-hand.depth} layers of width
\V{run.transformer-medium-hand.dim}, for a total of
\V{run.transformer-large-hand.parameters.millions} million parameters.
Compared with \run{transformer-medium-hand} it is
\V{pair.transformer-large-hand.transformer-medium-hand.verdict} on the test
set (gain \V{pair.transformer-large-hand.transformer-medium-hand.gain}; 95\%
CI \V{pair.transformer-large-hand.transformer-medium-hand.gain.low} to
\V{pair.transformer-large-hand.transformer-medium-hand.gain.high}), and the
modification is
\V{pair.transformer-large-hand.transformer-medium-hand.kept}.
\fi
\ifchangeattacks
\item \textbf{Attack-aware attention bias.} Compared with
\run{transformer-medium-hand}, \run{transformer-medium-hand-attacks} is
\V{pair.transformer-medium-hand-attacks.transformer-medium-hand.verdict} on
the test set (gain
\V{pair.transformer-medium-hand-attacks.transformer-medium-hand.gain}; 95\%
CI \V{pair.transformer-medium-hand-attacks.transformer-medium-hand.gain.low}
to
\V{pair.transformer-medium-hand-attacks.transformer-medium-hand.gain.high}),
and the modification is
\V{pair.transformer-medium-hand-attacks.transformer-medium-hand.kept}.
Compared with \run{transformer-medium}, which has no hand-crafted features,
\run{transformer-medium-attacks} is
\V{pair.transformer-medium-attacks.transformer-medium.verdict} on the test
set (gain \V{pair.transformer-medium-attacks.transformer-medium.gain}; 95\% CI
\V{pair.transformer-medium-attacks.transformer-medium.gain.low} to
\V{pair.transformer-medium-attacks.transformer-medium.gain.high}), and the
modification is
\V{pair.transformer-medium-attacks.transformer-medium.kept}.
\fi
\ifchangegeometry
\item \textbf{No relative position bias.} Compared with
\run{transformer-medium-hand},
\run{transformer-medium-hand-no-geometry} is
\V{pair.transformer-medium-hand-no-geometry.transformer-medium-hand.verdict}
on the test set (gain
\V{pair.transformer-medium-hand-no-geometry.transformer-medium-hand.gain};
95\% CI
\V{pair.transformer-medium-hand-no-geometry.transformer-medium-hand.gain.low}
to
\V{pair.transformer-medium-hand-no-geometry.transformer-medium-hand.gain.high}).
\fi
\end{itemize}

\ifchangeattacks
\paragraph{The attack-aware bias is the largest gain from random
initialisation.}
Without hand-crafted features it gains
\V{pair.transformer-medium-attacks.transformer-medium.gain} RMSE points,
compared with \V{pair.transformer-medium-long.transformer-medium.gain} for
the longer schedule and
\V{pair.transformer-medium-pretrained.transformer-medium.gain} for
pretraining, and the three intervals in Table~\ref{tab:pairs} do not overlap.
\run{transformer-medium-attacks}, trained from random initialisation, has a
test RMSE of \V{run.transformer-medium-attacks.score}, lower than the
\V{run.transformer-medium-pretrained.score} of the pretrained
\run{transformer-medium-pretrained}, and the intervals of their gains over
the GBDT baseline do not overlap (Table~\ref{tab:leaderboard}). With
hand-crafted features the bias gains
\V{pair.transformer-medium-hand-attacks.transformer-medium-hand.gain}, more
than the longer schedule
(\V{pair.transformer-medium-hand-long.transformer-medium-hand.gain}) and less
than pretraining
(\V{pair.transformer-medium-hand-pretrained.transformer-medium-hand.gain}).

\paragraph{With the attack-aware bias, hand-crafted features add almost
nothing.}
The test RMSE of \run{transformer-medium-hand-attacks} is
\V{run.transformer-medium-hand-attacks.score} and that of
\run{transformer-medium-attacks} is
\V{run.transformer-medium-attacks.score}, a difference of
\V{run.transformer-medium-hand-attacks.over-twin} points, where the paired
gain from the features is
\V{pair.transformer-medium-hand.transformer-medium.gain} without the bias.
The two attack-aware models were not compared directly on resampled games,
and their gains over the GBDT baseline lie within each other's intervals
(Table~\ref{tab:leaderboard}). The result agrees with that of the longer schedule: the
transformer can derive most of what the hand-crafted features encode from the
board. We hypothesise that, given the attack relations, it does so within
the standard number of epochs.

\paragraph{Both attack-aware models were still improving when training
ended.}
Both had their best validation score in the final epoch
(\V{run.transformer-medium-attacks.best-epoch} of
\V{run.transformer-medium-attacks.epochs}; Figure~\ref{fig:shapes}). Neither
is the best model of this task, which remains \V{best.name} with an RMSE of
\V{best.score}. The two modifications can be combined: the additional biases
are initialised to zero, so an attack-aware transformer can start from the
same pretrained encoder and initially computes the same function. This
combination is the next experiment for this task.
\fi

\paragraph{Changes that did not help.}
\ifchangelarger
The larger transformer required \V{run.transformer-large-hand.seconds} of
GPU time, compared with \V{run.transformer-medium-hand.seconds} for the
medium transformer, and is less accurate. Its validation curve lies above
that of the medium transformer after every epoch and flattens without
rising, so the model is not overfitting (Figure~\ref{fig:shapes}). We
hypothesise that it learns more slowly at the medium transformer's learning
rate and number of epochs. Whether a learning rate or schedule tuned for the
larger model would change the result was not tested.
\fi
\ifchangehybrid
The hybrid shows the opposite pattern. It is ahead of the transformer for
most of the schedule, has its best validation score at epoch
\V{run.hybrid-medium-hand.best-epoch} of \V{run.hybrid-medium-hand.epochs},
and then deteriorates while the transformer continues to improve.
\fi
\ifchangeconv
The ResNet behaves like the hybrid but starts from a larger lead, with its
best epoch at \V{run.conv-medium-hand.best-epoch}, and the epoch selected on
the validation set retains a small gain. We hypothesise that the models with
convolutions fit the training puzzles faster and overfit sooner, so that
they would benefit more from additional data than from additional epochs. The
ResNet is also the cheapest of the medium models to train
(\V{run.conv-medium-hand.seconds}). Because its learning rate differs from
the transformer's, its gain cannot be attributed to the architecture alone.
\fi
\ifchangegeometry
Removing the relative position bias costs less than the attack-aware bias
adds.
\fi
Each of these models was trained once.
\fi

% Training-set size: the same ablation for every task, so a model paper
% inputs this file inside its ablation section and then reports its own
% findings. It appears once the task has the runs.
% The prose follows docs/writing.md.
\ifgames
\subsection{Training-set size}
\label{sec:training-size}

To estimate how much a model's accuracy depends on the amount of training
data, we retrain it on a half and on a quarter of the training games and
evaluate it on the full test set (Table~\ref{tab:games},
Figure~\ref{fig:games}). The subsets are selected by a hash of the game
identifier and are nested, so the quarter is contained in the half. A network
is given twice the epochs on half the games and four times the epochs on a
quarter, so the number of optimisation steps is constant and only the amount
of data varies.
\ifmore
We also train on more data. A larger sample of the same source is split in
the same way as the standard dataset, so the validation and test sets are
unchanged and no test game appears in the larger training set. A network
trained on a larger set is given proportionally fewer epochs, so that it
again takes the same number of optimisation steps as its base. Unlike a model
trained on a subset, a model trained on more data is a candidate model, and
the acceptance criterion applies to it.
\fi

\begin{table}[htbp]\centering\footnotesize
\block{study-games}
\caption{Training-set size. Each model is trained on a fraction of the
training games\ifmore, or on a multiple of them from the larger sample,\fi{}
and evaluated on the full test set. The gain is paired with the same model
trained on the standard training set. A negative gain means that the model is
worse, and bold marks statistical significance.}
\label{tab:games}
\end{table}

\begin{figure}[htbp]\centering
\block{study-games-figure}
\caption{Test score as a function of the number of training games. Each step
to the right doubles the training set.}
\label{fig:games}
\end{figure}
\fi


\ifgames
\paragraph{Smaller training sets.}
Both model classes are less accurate with less data, and the transformer
loses more. Relative to the same model trained on the full training set,
\run{transformer-medium-hand} has a paired gain of
\V{pair.transformer-medium-hand-half.transformer-medium-hand.gain} RMSE
points on half of the training puzzles and of
\V{pair.transformer-medium-hand-quarter.transformer-medium-hand.gain} on a
quarter. For GBDT the gains are \V{pair.gbm-hand-half.gbm-hand.gain} and
\V{pair.gbm-hand-quarter.gbm-hand.gain}. From a quarter of the puzzles to
all of them, the transformer's advantage over GBDT grows with every doubling
(Table~\ref{tab:games}), and its curve in Figure~\ref{fig:games} is still
falling at the full training set.
\ifchangeepochs
The last of these doublings is worth about as much to the transformer as
doubling the number of epochs
(\V{pair.transformer-medium-hand-long.transformer-medium-hand.gain}).
\fi
The dataset is a \V{recipe.sample.percent}\% sample of the puzzles with a
stable rating, so more data is available. From these results alone one would
expect additional puzzles to help the transformer more than a longer
schedule\ifmore. The larger training sets do not confirm this
expectation\fi.
\fi

\ifmore
\paragraph{Larger training sets.}
The larger sample contains up to \V{more.gbm-hand.most} the training puzzles,
with the validation and test sets unchanged. Relative to the same model on
the standard training set, \run{transformer-medium-hand} has a paired gain of
\V{pair.transformer-medium-hand-double.transformer-medium-hand.gain} RMSE
points on twice the training puzzles
(\V{pair.transformer-medium-hand-double.transformer-medium-hand.verdict};
\V{pair.transformer-medium-hand-double.transformer-medium-hand.kept}) and of
\V{pair.transformer-medium-hand-quadruple.transformer-medium-hand.gain} on
four times
(\V{pair.transformer-medium-hand-quadruple.transformer-medium-hand.verdict};
\V{pair.transformer-medium-hand-quadruple.transformer-medium-hand.kept}).
For GBDT the gains are \V{pair.gbm-hand-double.gbm-hand.gain}
(\V{pair.gbm-hand-double.gbm-hand.verdict};
\V{pair.gbm-hand-double.gbm-hand.kept}) and
\V{pair.gbm-hand-quadruple.gbm-hand.gain}
(\V{pair.gbm-hand-quadruple.gbm-hand.verdict};
\V{pair.gbm-hand-quadruple.gbm-hand.kept}).

More data helps both model classes, and it helps GBDT more: on four times
the puzzles, the intervals of the two paired gains do not overlap
(Table~\ref{tab:games}). For the transformer, four times the puzzles yield a
smaller gain than twice the puzzles. The intervals of these two gains do not
overlap either, although the two models were not compared directly and each
was trained once. No model trained on a larger set approaches the best
model, which remains \V{best.name}.

The three transformer models take the same number of optimisation steps, and
each had its best validation score in its final epoch:
\V{run.transformer-medium-hand.best-epoch} of
\V{run.transformer-medium-hand.epochs} on the standard training set,
\V{run.transformer-medium-hand-double.best-epoch} of
\V{run.transformer-medium-hand-double.epochs} on twice the puzzles and
\V{run.transformer-medium-hand-quadruple.best-epoch} of
\V{run.transformer-medium-hand-quadruple.epochs} on four times. On half and
on a quarter of the puzzles, where each puzzle is seen two and four times as
often, the best epoch was \V{run.transformer-medium-hand-half.best-epoch} of
\V{run.transformer-medium-hand-half.epochs} and
\V{run.transformer-medium-hand-quarter.best-epoch} of
\V{run.transformer-medium-hand-quarter.epochs}. We hypothesise that part of
the transformer's loss on the smaller training sets is due to overfitting,
which data beyond the standard training set cannot recover, and that from
the standard training set upward the transformer is limited by its number of
optimisation steps.
\ifchangeepochs
Consistent with this, doubling the epochs on the standard training set gains
\V{pair.transformer-medium-hand-long.transformer-medium-hand.gain} RMSE
points, more than either larger training set (the intervals do not overlap).
\fi

GBDT was constrained as well. A model may use at most
\V{method.gbm.rounds} boosting rounds. On four times the puzzles it used
\V{run.gbm-hand-quadruple.trees} and on twice the puzzles
\V{run.gbm-hand-double.trees}, whereas on the standard training set it
stopped at \V{run.gbm-hand.trees}. The gain of GBDT from four times the data
is therefore a lower bound on what that data allows.
\fi

% Feature-group ablation: the same ablation for every task, so a model paper
% inputs this file inside its ablation section and then reports its own
% findings. It appears once the task has the runs.
% The prose follows docs/writing.md.
\iffeatures
\subsection{Feature groups}
\label{sec:feature-groups}

\ifsequence
The hand-crafted features fall into three groups: move accuracy, clock use
and the number of available moves.
\else
The hand-crafted features fall into groups: material, mobility, hanging
pieces, passed pawns and king safety, and, where the task gives a line of
moves, the opponent's last move, the first move of the line and the line as a
whole.
\fi
We retrain a model with one group removed at a time
(Table~\ref{tab:features}, Figure~\ref{fig:features}). A significant loss
shows that the model used the group and that no other input provides the same
information. The absence of a loss has two possible explanations: the group
is uninformative, or the same information is available from another group or
from the other inputs. The ablation therefore measures the marginal
contribution of a group given everything else the model receives, which can
be much smaller than the information the group carries on its own.

\begin{table}[htbp]\centering\footnotesize
\block{study-features}
\caption{Feature-group ablation. Each model is trained on all hand-crafted
features except one group. The gain is paired with the model trained on all
features. A negative gain means that the model is worse without the group,
and bold marks statistical significance. Groups are ordered by the size of
the loss.}
\label{tab:features}
\end{table}

\begin{figure}[htbp]\centering
\block{study-features-figure}
\caption{Paired gain of the GBDT baseline without one feature group over the
GBDT baseline with all features, with 95\% CIs. Left of the vertical line,
the model is worse without the group.}
\label{fig:features}
\end{figure}
\fi


\iffeatures
\paragraph{The solution line is the most important group.}
The GBDT baseline is significantly worse in
\V{study.features.gbm-hand.worse.of}. The largest loss follows the removal of
the \V{study.features.gbm-hand.most.group} group, with a paired gain of
\V{study.features.gbm-hand.most.gain} RMSE points, even though the solution
length remains available as a context feature. Next are mobility
(\V{pair.gbm-hand-no-mobility.gbm-hand.gain}) and the first move of the
solution (\V{pair.gbm-hand-no-first-move.gbm-hand.gain}). The remaining groups
matter little: material (\V{pair.gbm-hand-no-material.gbm-hand.gain}), king
safety (\V{pair.gbm-hand-no-kings.gbm-hand.gain}), hanging pieces
(\V{pair.gbm-hand-no-hanging.gbm-hand.gain}) and passed pawns
(\V{pair.gbm-hand-no-pawns.gbm-hand.gain}). Without the single feature that
describes the opponent's last move, the baseline is
\V{study.features.gbm-hand.least.verdict}, with a gain of
\V{study.features.gbm-hand.least.gain}. We do not interpret a difference of
this size, since every model was trained with a single seed. The GBDT
baseline's information about a puzzle's difficulty thus comes mainly from the
properties of its solution. The companion task that withholds the solution
measures the same dependence directly, and its paper reports the result.
\fi
\fi

\section{Error Analysis}
\label{sec:error}

% Error analysis: what the analysis is and how it is computed, the same for
% every task. A model paper inputs this file at the start of its error
% analysis, before the subgroups that are specific to its task.
% The prose follows docs/writing.md.
\iflook
The main results identify the best model. They do not show where a model is
accurate and where it is not, which determines how it can be used. We
therefore analyse the test predictions of every model by subgroup. All
results in this section are computed from the \V{look.test-rows} test
examples (analysis of \V{look.date}).

\paragraph{Subgroups.}
We partition the test set by one variable at a time. For each subgroup we
report the number of test examples, the mean observed target, the mean
prediction of a model and the model's score. Two models are compared:
\V{challenger.name}, the best model other than the GBDT baseline, and the
GBDT baseline itself. The gain of the first over the second is computed
within each subgroup as for the whole test set, as a paired difference over
resampled games (\V{look.resamples} resamples), and is omitted for subgroups
with fewer than \V{look.smallest-compared} test examples. A subgroup whose
interval contains zero provides no evidence in either direction. A numeric
variable is divided into ranges, and a range written $a$-$b$ includes $a$ and
excludes $b$.

\paragraph{Two-way tables.}
For a pair of variables we show the mean observed target and the mean
prediction in each cell, on a common colour scale. Cells with fewer than
\V{grid.smallest-cell} test examples are left empty.

\paragraph{Examples.}
Example \ifsequence games\else positions\fi{} are selected by fixed rules and
not by inspection: typical examples chosen by a hash of their key, the
examples with the best model's most extreme predictions in each direction,
and its largest errors in each direction, \V{examples.count} of each. They
illustrate the quantities under discussion and are not evidence of how often
a pattern occurs.
\else
\emph{This version of the paper was built before the error analysis was run.
The analysis partitions the test set by rating, clock and properties of the
position, measures calibration, and shows example positions. It is produced
by \texttt{lab analyse} on the machine that holds the dataset and the test
predictions of every model, and this section appears once its output is
committed.}
\fi


\iflook
\subsection{Distribution of the errors}

The median absolute error of the best model is \V{best.close.median} rating
points, and 90\% of its predictions are within \V{best.close.p90} points of
the true rating. \V{best.close.within.100}\% of the predictions are within
100 points, \V{best.close.within.200}\% within 200 points and
\V{best.close.within.300}\% within 300 points (Table~\ref{tab:closeness}).
With rating classes 200 points wide, the first figure is the share of
puzzles placed within half a class of their true rating, and the second the
share placed within a whole class.

\begin{table}[htbp]\centering\small
\block{closeness}
\caption{Distribution of the absolute errors on the test set: the median,
the 90th percentile, and the share of predictions within 100, 200 and 300
rating points of the true rating.}
\label{tab:closeness}
\end{table}

\subsection{Calibration}

The best model's predictions are accurate on average at every level. In
Figure~\ref{fig:reliability} the test puzzles are divided into
\V{best.reliability.groups} bins of equal size by the model's prediction. In
the lowest bin the mean prediction is \V{best.reliability.lowest.said} and
the mean true rating is \V{best.reliability.lowest.happened}. In the highest
bin they are \V{best.reliability.highest.said} and
\V{best.reliability.highest.happened}. The largest difference in any bin is
\V{best.reliability.widest-gap} points. The deviations have a consistent
direction: the lowest predictions are slightly too low and the highest
slightly too high, so the predictions are somewhat more spread out than the
true ratings.

Predictions for individual puzzles are much less precise. For a given
prediction, the central 80\% of the true ratings span about
\V{best.reliability.spread} points (shaded band in
Figure~\ref{fig:reliability}).

\begin{figure}[htbp]\centering
\block{reliability}
\caption{Mean true rating against mean predicted rating. The test puzzles
are divided into bins of equal size by the prediction, and each point is one
bin. The shaded band is the central 80\% of the true ratings in each bin of
the first model. Points on the diagonal are unbiased.}
\label{fig:reliability}
\end{figure}

Conditioning on the true rating instead of the prediction shows regression
towards the mean (Figure~\ref{fig:band-bias}). For puzzles with a true rating
below 1000, the mean prediction of \V{challenger.name} minus the mean true
rating is \V{slice.band.-1000.challenger.off} points, and for puzzles rated
2200 and above it is \V{slice.band.2200+.challenger.off} points: easy
puzzles are rated too hard and hard puzzles too easy. The GBDT baseline is more biased in both
bands (\V{slice.band.-1000.reference.off} and
\V{slice.band.2200+.reference.off}).

The two analyses answer different questions. When a predicted rating is
shown for a tactic, the first applies: tactics predicted to be rated 1800
are rated 1800 on average. When the question is whether the model recognises
a truly hard tactic, the second applies: the model rates it as hard, but
usually not as hard as it is. Every model with prediction error shows this
effect, because examples it cannot distinguish receive predictions near
their common mean, and a more accurate model shows less of it.

\begin{figure}[htbp]\centering
\block{slice-band-bias}
\caption{Mean prediction minus mean true rating, by band of the true rating,
for the two models compared. Above the line a model rates the puzzles of the
band as harder than they are, and below it as easier.}
\label{fig:band-bias}
\end{figure}

\subsection{Solution length}

Table~\ref{tab:moves} gives the results by the number of moves the solver has
to find. \V{challenger.name} is significantly better than the GBDT baseline
in \V{slice.moves.ahead.of}. Its gain is smallest for one-move puzzles
(\V{slice.moves.1.gain} RMSE points) and larger for every longer solution
(\V{slice.moves.2.gain} at two moves and \V{slice.moves.3.gain} at three). A
plausible explanation is that a one-move solution has few properties beyond
those the hand-crafted features already state, whereas a longer solution
requires foreseeing the opponent's replies, for which the board carries
information that the features do not.

\begin{table}[htbp]\centering\footnotesize
\block{slice-moves}
\caption{Test-set results by solution length. For each group: the number of
test puzzles, the mean true rating, the mean prediction and RMSE of each of
the two models, and the paired gain of the first model over the second.}
\label{tab:moves}
\end{table}

The GBDT baseline is nearly unbiased at every solution length, as expected of
a model that can split directly on this feature. The network is less well
calibrated by length: for one-move puzzles its mean prediction minus the mean
true rating is \V{slice.moves.1.challenger.off} points
(Figure~\ref{fig:slice-moves}).

\begin{figure}[htbp]\centering
\block{slice-moves-bias}
\caption{Mean prediction minus mean true rating, by solution length, for the
two models compared.}
\label{fig:slice-moves}
\end{figure}

\subsection{Themes}

Lichess tags every puzzle with themes that describe the tactical motif (for
example a fork, a pin or a discovered attack), the material won, the phase of
the game and the length of the solution. No model receives the themes as
input. We form one group for each of the \V{slice.themes.compared} most
common themes, and a puzzle belongs to the group of every theme it has.
\V{challenger.name} is significantly better than the GBDT baseline in
\V{slice.themes.ahead.of} (Figure~\ref{fig:themes}).

The gain is largest for the theme
\texttt{\V{slice.themes.gain.most.group}}
(\V{slice.themes.gain.most} RMSE points) and smallest for
\texttt{\V{slice.themes.gain.least.group}} (\V{slice.themes.gain.least}).
After pawn endgames, the largest gains are for pins
(\V{slice.themes.pin.gain}) and quiet moves
(\V{slice.themes.quietMove.gain}). After promotions, the smallest are for
hanging pieces (\V{slice.themes.hangingPiece.gain}) and one-move puzzles
(\V{slice.themes.oneMove.gain}). We hypothesise that the network adds most
where the difficulty depends on the configuration of the pieces, which the
hand-crafted features summarise only coarsely: a pin is a relation between
three pieces on a line, and a pawn endgame is often decided by a single
tempo. In a promotion or the capture of a hanging piece, a few features
describe most of the tactic. Table~\ref{tab:slice-themes} in the appendix
gives the values behind the figure.

\begin{figure}[htbp]\centering
\block{slice-themes-gain}
\caption{Paired gain in RMSE of the best model other than the GBDT baseline
over the GBDT baseline, for the \V{slice.themes.compared} most common themes.
A puzzle has several themes and is counted under each. A theme whose interval
includes zero provides no evidence in either direction.}
\label{fig:themes}
\end{figure}

\subsection{Quiet moves and mates}

The two most important features of the GBDT baseline correspond to large
differences in the mean rating. A quiet move is a move that is neither a
capture nor a check, which the solver has to consider without a forcing cue
(Table~\ref{tab:quiet}). Puzzles whose solution has no quiet move have a mean
rating of \V{slice.quiet.0.happened}, and puzzles with one quiet move a mean
rating of \V{slice.quiet.1.happened}. A second quiet move adds much less
(\V{slice.quiet.2.happened}): most of the difficulty lies in the presence of
a quiet move at all.

A solution that ends in mate is forcing throughout
(Table~\ref{tab:mate}). Such puzzles have a mean rating of
\V{slice.mate.yes.happened}, compared with \V{slice.mate.no.happened} for
puzzles whose solution wins material.

The network is significantly better than the GBDT baseline in every group of
both tables. Its gain is \V{slice.quiet.2.gain} RMSE points for solutions
with two quiet moves and \V{slice.quiet.0.gain} for solutions with none, and
the two intervals do not overlap. Its gains for solutions that end in mate
(\V{slice.mate.yes.gain}) and for those that do not (\V{slice.mate.no.gain})
are too close to distinguish.

\begin{table}[htbp]\centering\footnotesize
\block{slice-quiet}
\caption{Test-set results by the number of quiet moves (neither captures nor
checks) among the solver's moves.}
\label{tab:quiet}
\end{table}

\begin{table}[htbp]\centering\footnotesize
\block{slice-mate}
\caption{Test-set results by whether the solution ends in mate.}
\label{tab:mate}
\end{table}

\subsection{Properties of the first move}

The first move of the solution is informative about the difficulty of a
puzzle. The mean true ratings below depend only on the puzzles and not on
any model.
\begin{itemize}
\item \textbf{Checks are easier to find.} Puzzles whose first move gives
check have a mean rating of \V{slice.check.yes.happened}, compared with
\V{slice.check.no.happened} for the others.
\item \textbf{Captures make little difference.} The mean rating is
\V{slice.capture.yes.happened} when the first move is a capture and
\V{slice.capture.no.happened} when it is not.
\item \textbf{Sacrifices are harder.} When the first move leaves the moved
piece open to capture, the mean rating is \V{slice.sacrifice.yes.happened},
compared with \V{slice.sacrifice.no.happened} otherwise.
\item \textbf{Retreats are harder.} When the first move goes backward, the
mean rating is \V{slice.backward.yes.happened}, compared with
\V{slice.backward.no.happened} otherwise.
\item \textbf{The piece matters.} Puzzles that begin with a king move have
the highest mean rating (\V{slice.first-piece.king.happened}), followed by
those that begin with a pawn move (\V{slice.first-piece.pawn.happened}).
Those that begin with a queen move have the lowest
(\V{slice.first-piece.queen.happened}).
\end{itemize}
The network is significantly better than the GBDT baseline in every one of
these groups (Tables~\ref{tab:slice-first-piece}
to~\ref{tab:slice-backward} in the appendix), so its gain is not specific to
any kind of first move.

\subsection{Reliability of the rating}

The network's gain over the GBDT baseline depends on how well a puzzle's
rating is measured. By rating deviation, the gain is
\V{slice.deviation.-75.gain} RMSE points for the most reliable ratings (a
deviation below 75) and \V{slice.deviation.90+.gain} for the least reliable
(a deviation of 90 or more), and it increases from each group to the next
(Table~\ref{tab:slice-deviation}). By number of attempts, the gain lies
between \V{slice.plays.-200.gain} and \V{slice.plays.200-500.gain} in the
three groups of puzzles attempted fewer than 1{,}000 times and falls to
\V{slice.plays.5000+.gain} for puzzles attempted 5{,}000 times or more
(Table~\ref{tab:slice-plays}).

The gain is therefore smallest for the puzzles whose ratings are measured
best and for the puzzles attempted most often. The subgroup tables do not
explain why. The two groups also differ in difficulty: the most reliable
ratings have a mean of \V{slice.deviation.-75.happened} and the least
reliable a mean of \V{slice.deviation.90+.happened}, and the gain varies with
the rating band as well (Table~\ref{tab:bands}). An analysis by rating
deviation within each rating band would separate the two effects and has not
been done. Until then, the conservative conclusion is that the network's
gain is larger on the less frequently attempted puzzles and small on the
puzzles whose ratings are measured best.

\subsection{Solution length and material}

Figure~\ref{fig:grid} shows the mean rating by solution length and by the
number of minor and major pieces on the board, which serves as a proxy for
the phase of the game. The agreement between the two panels indicates that
the model reproduces the joint pattern and not only the two marginal
effects.

\begin{figure}[htbp]\centering
\block{grid-moves-pieces}

\medskip
\block{grid-moves-pieces-said}
\caption{Mean rating by solution length and by the number of knights,
bishops, rooks and queens on the board. Top: true ratings. Bottom:
predictions of the best model other than the GBDT baseline. Both panels use
the same colour scale, and cells with too few test puzzles are left empty.}
\label{fig:grid}
\end{figure}

\subsection{Agreement with the GBDT baseline}

Two models with similar scores can still disagree on individual puzzles, in
which case an average of their predictions may be more accurate than either.
Table~\ref{tab:beside} gives, for each model, the correlation of its
predictions with those of the GBDT baseline and the share of test puzzles on
which it is the closer of the two to the true rating. The best model is
closer than the GBDT baseline on \V{best.beside.closer}\% of the test
puzzles, and the correlation between the two sets of predictions is
\V{best.beside.correlation}. The two models therefore still disagree on many
individual puzzles, which suggests that an ensemble of the two could improve
on both. We have not trained one.

\begin{table}[htbp]\centering\small
\block{beside}
\caption{Agreement with the GBDT baseline on the test set: the correlation
between the two models' predictions, and the share of puzzles on which the
model is closer to the true rating than the baseline.}
\label{tab:beside}
\end{table}

\subsection{Example positions}
\label{sec:positions}

The following examples were selected from the test set by the fixed rules
described above. Each board is shown from the mover's side. The grey arrow is
the opponent's preceding move and the blue arrow is the first move of the
solution. Next to each board we give the solution, the true rating, the
predictions of the best model and of the GBDT baseline, and the puzzle's
identifier on Lichess.

The largest errors are the most informative, because they show what the
inputs do not contain. Most of the puzzles whose rating the model
overestimates or underestimates by the most carry the theme
\texttt{endgame} and have few pieces on the board (Figures~\ref{fig:over}
and~\ref{fig:under}). This agrees with the subgroup results: the model's
RMSE is highest for positions with fewer than three minor and major pieces
(\V{slice.pieces.-3.challenger}, compared with
\V{slice.pieces.12+.challenger} for positions with twelve or more;
Table~\ref{tab:slice-pieces}). We hypothesise that the difficulty of such
positions depends on calculation that neither the board encoding nor the
hand-crafted features summarise. The puzzles with the highest and the lowest
predictions are shown in the appendix (Figures~\ref{fig:highest}
and~\ref{fig:lowest}).

\examples{examples-typical}{Typical puzzles: the first \V{examples.count}
test puzzles in the order of a hash of their key.}{fig:typical}

\examples{examples-over}{The largest overestimates: the test puzzles for
which the best model's prediction exceeds the true rating by the most.}{fig:over}

\examples{examples-under}{The largest underestimates: the test puzzles for
which the best model's prediction falls below the true rating by the
most.}{fig:under}
\fi

\section{Practical Use}
\label{sec:use}

\begin{itemize}
\item \textbf{Ordering the tactics of a game.} With a Spearman correlation
of \V{best.spearman}, the model can identify which of a game's missed
tactics was routine and which was hard. The results support this use.
\item \textbf{Comparing a tactic with a player's level, with a margin.}
Compared with a player's own puzzle rating, a predicted rating indicates
whether a missed tactic was within the player's reach. Because the MAE is
\V{best.mae} rating points, such a comparison should leave a margin of about
one rating class on either side and distinguish only tactics well below the
player's level, well above it, or near it.
\item \textbf{Not for display as a number.} A predicted rating is accurate on
average and imprecise for a single tactic, and it underestimates the hardest
tactics. Displayed as a number, it would suggest a precision that the model
does not have.
\item \textbf{GBDT provides most of the accuracy.} The GBDT baseline has an
RMSE of \V{reference.score}, compared with \V{best.score} for the best
model, and needs no GPU. Where the application has to run the model itself,
GBDT is the natural starting point, and the network is an improvement that a
server could add.
\end{itemize}

\section{Limitations}
\label{sec:limits}

\begin{itemize}
\item \textbf{The target is on the Lichess puzzle scale.} Relating a
prediction to a player's game rating, or to another site's scale, requires a
mapping that this work does not provide.
\item \textbf{The models are given the solution.} How hard a tactic is to
notice when no solution is given is a different question. The companion task
that withholds the solution addresses it on the same puzzles, and its paper
compares the two.
\item \textbf{Puzzle ratings depend on who attempts the puzzle.} Lichess
presents puzzles near a solver's level, so a hard puzzle is rated mostly by
strong solvers and an easy one mostly by weaker solvers. We take the rating
as given.
\item \textbf{Only stable ratings are included.} Puzzles with few attempts
are excluded, so the results do not show how the models perform on puzzles
that have rarely been attempted, including new ones.
\item \textbf{Every position contains a winning tactic.} A puzzle is taken
from a position in which a winning line exists. The models have seen no
position in which a tempting tactic fails, and their predictions for such
positions are untested.
\item \textbf{Single training runs.} Each network was trained with one
random seed. The confidence intervals reflect the sampling of test games and
not the variation between training runs.
\end{itemize}

\section{Reproducibility}
\label{sec:cost}

Training the \V{runs.count} models on the leaderboard took \V{cost.seconds}
in total: \V{cost.seconds.networks} for the networks, each trained on a
single cloud GPU, and \V{cost.seconds.others} for the constant baseline and
GBDT on a CPU. The ablation models that are not on the leaderboard are not
included in this total. The best model has \V{best.parameters} parameters and
was trained for \V{best.epochs.run} epochs in \V{best.seconds} on an
\V{best.device}, at about \V{best.rows-a-second} puzzles per second, not
counting the pretraining of its encoder. Table~\ref{tab:runs} lists
the hardware and training time of every model, and
Table~\ref{tab:settings} the hyperparameters of the networks of the fixed
sequence.

All results are generated from committed files. The dataset is built from
its recipe by \texttt{lab data puzzles}, the models are trained by
\texttt{lab all puzzle-difficulty}, each training run writes a record to the
model's \texttt{results/} directory, and
\texttt{lab paper puzzle-difficulty} generates every number, table and
figure of this paper from those records and compiles it. A record names the
build of the dataset it used (\V{data.fingerprint}) and the commit of the
code that produced it (\V{code.commit}). Model weights and test predictions
are not stored in the repository. The model's card states where they are
kept and gives the hash of the best model's weights.

\section{Conclusion and Future Work}
\label{sec:conclusion}

The rating of a Lichess puzzle can be predicted from its position and
solution with an RMSE of \V{best.score} rating points, compared with
\V{floor.score} for a constant prediction. GBDT on hand-crafted features of
the solution accounts for most of this improvement, and a transformer that
reads the board adds a significant further gain, which is largest for the
easiest and the hardest puzzles. The predictions order puzzles reliably but
are too imprecise to report as individual ratings. The following experiments
are the natural next steps.
\begin{itemize}
\ifchangeepochs\else
\item \textbf{A longer schedule.} Most networks were still improving when
training ended, so additional epochs are the cheapest experiment to add.
\fi
\iffeatures\else
\item \textbf{Feature-group ablation.} Removing one group of hand-crafted
features at a time would show which of them a model depends on.
\fi
\ifchangeattacks\ifchangestart
\item \textbf{Attack-aware attention bias with pretraining.} Adding the
attack-aware bias to \run{transformer-medium-hand-pretrained} changes one
factor in the best model and combines the two modifications from which the
transformer gained most.
\fi\fi
\ifmore
\item \textbf{More data together with more optimisation steps.} On the
larger training sets the transformer was still improving when its schedule
ended. Training on four times the puzzles with twice the steps, and
training the best model on four times the puzzles, are the two missing
cases.
\item \textbf{A higher limit on boosting rounds.} On four times the puzzles
GBDT used every round it was allowed.
\fi
\ifchangestart
\item \textbf{A longer schedule from the pretrained encoder.}
\run{transformer-medium-hand-pretrained} had its best validation score at
epoch \V{run.transformer-medium-hand-pretrained.best-epoch} of
\V{run.transformer-medium-hand-pretrained.epochs}, so it was stopped by its
schedule and not by early stopping.
\item \textbf{A stronger pretrained encoder.} The encoder was pretrained
once, at one size, on part of the evaluation database. Longer pretraining,
more positions and a larger encoder have not been tried.
\fi
\item \textbf{An encoder pretrained on human moves.} The winning entry of
the 2025 competition fine-tunes a human-move
model~\cite{puzzlecupwinner}. Using the Maia-2 encoder under this task's
prediction head is a natural comparison.
\item \textbf{An ensemble.} The best network and the GBDT baseline disagree
often enough that an average of the two may be more accurate than either, at
the cost of running both.
\item \textbf{Calibration to a player's scale.} Presenting a prediction to a
player requires a mapping to the player's own rating scale, and validation
on tactics from real games.
\end{itemize}

\bibliographystyle{unsrt}
\bibliography{../../../papers/references}

\clearpage
\appendix

\section{Hyperparameters and Features}

\begin{table}[htbp]\centering\footnotesize
\block{settings}
\caption{Hyperparameters of the networks of the fixed sequence, from their
configuration files. \emph{LR} is the peak learning rate. Values shared by
all of them are given below the table.}
\label{tab:settings}
\end{table}

\begin{table}[htbp]\centering\footnotesize
\block{hand}
\caption{Hand-crafted features: the inputs of \V{task.reference} in addition
to the scalar context features, and of a network whose configuration includes
them.}
\label{tab:hand}
\end{table}

\iflook
\section{Examples with Extreme Predictions}

\examples{examples-highest}{The test puzzles with the highest predictions of
the best model.}{fig:highest}

\examples{examples-lowest}{The test puzzles with the lowest predictions of
the best model.}{fig:lowest}

\section{Subgroup Tables}

Each table below partitions the test puzzles by one variable. The columns are
the group, the number of test puzzles, the mean true rating, the mean
prediction and RMSE of each of the two models compared, and the paired gain
of the first model over the second with its 95\% CI, in bold when
statistically significant.

\block{slices}
\fi

\end{document}
